\chapter{How to Help Someone Who Is Misusing Drugs or Alcohol}

\section*{Definition and Overview}
Helping someone who is misusing drugs or alcohol can feel overwhelming, but family and friends play a vital role in recovery. Effective help starts with compassion, honest conversation, and practical support rather than judgement or ultimatums. The goals are to encourage the person to seek help, connect them with treatment services, and support them through change while looking after your own wellbeing. Recovery is often a long process with relapses, and the most helpful responses are patient, informed, and consistent.

\section*{Epidemiology}
Around one in ten Australians drinks at levels that put their health at risk, and a significant number use illicit drugs. Alcohol and other drug problems affect people of all ages and backgrounds, and only a minority ever receive formal treatment. Family members often seek help first, and family involvement is associated with better treatment engagement and outcomes. Support services for families, such as counselling and helplines, are available across Australia.

\section*{Aetiology and Pathophysiology}
\begin{itemize}
    \item \textbf{Why people use:} pleasure, coping with stress or trauma, social pressure, and physical dependence
    \item \textbf{Dependence:} the brain adapts to the substance, so stopping causes cravings and withdrawal
    \item \textbf{Harm:} health, financial, legal, and relationship consequences accumulate
    \item \textbf{Ambivalence:} most people are genuinely conflicted about their use
    \item \textbf{Change process:} recovery usually involves stages from contemplation through action to maintenance, with relapses common
    \item \textbf{Family dynamics:} enabling, conflict, and codependency can unintentionally sustain use
    \item \textbf{Pathophysiology of recovery:} treatment addresses both the physical dependence and the psychological and social drivers
\end{itemize}

\section*{Clinical Features}
\begin{itemize}
    \item Changes in behaviour, mood, and reliability
    \item Withdrawal from family and usual activities
    \item Financial problems and borrowing money
    \item Declining work, study, or health
    \item Denial, secrecy, or defensiveness about use
    \item Signs of intoxication or withdrawal
    \item The family's own stress, worry, and burnout
\end{itemize}

\section*{Differential Diagnosis}
\begin{itemize}
    \item Mental health conditions such as depression and anxiety can coexist with and drive substance use
    \item Physical symptoms may reflect the substance, withdrawal, or other illness
    \item The type of substance affects risks and treatment
    \item Dependence, harmful use, and recreational use need different responses
    \item A professional assessment clarifies the situation
\end{itemize}

\section*{When to Seek Medical Care}
Seek help from a GP, drug and alcohol service, or helpline when concerned about someone's use. Encouraging the person to have a health assessment is often the first step. Seek urgent care for serious physical or mental health problems.

\textbf{Call 000 (or local emergency services) for:}
\begin{itemize}
    \item Unconsciousness or difficulty waking the person
    \item Seizures, severe confusion, or hallucinations
    \item Difficulty breathing or chest pain
    \item Signs of overdose, including slow breathing and pinpoint pupils
    \item Suicidal thoughts, self-harm, or violence
    \item Severe withdrawal with agitation, fever, or vomiting
\end{itemize}

\section*{Diagnosis and Investigations}
\begin{itemize}
    \item \textbf{Observation:} noting patterns, consequences, and risks without judgement
    \item \textbf{Conversation:} talking calmly when the person is sober and receptive
    \item \textbf{Professional assessment:} a GP or drug and alcohol service assesses the pattern and severity of use
    \item \textbf{Screening:} for mental health conditions and physical harm
    \item \textbf{Family support assessment:} services assess and support family members
    \item \textbf{Planning:} matching treatment to the person's needs and readiness
\end{itemize}

\section*{Management}
\begin{itemize}
    \item \textbf{Start with compassion:} express concern for their wellbeing, not criticism of their behaviour
    \item \textbf{Talk at the right time:} choose calm moments when they are sober
    \item \textbf{Listen:} understand their perspective and reasons for using
    \item \textbf{Encourage professional help:} offer to accompany them to a GP or treatment service
    \item \textbf{Support treatment:} attend appointments, help with logistics, and celebrate progress
    \item \textbf{Set boundaries:} avoid enabling behaviours such as providing money for substances
    \item \textbf{Plan for crises:} know emergency numbers and overdose responses
    \item \textbf{Look after yourself:} seek counselling and support for your own stress
\end{itemize}

\section*{Complications}
\begin{itemize}
    \item Overdose, injury, and other medical emergencies
    \item Withdrawal complications, including seizures
    \item Mental health deterioration and suicide risk
    \item Relationship breakdown, family conflict, and domestic violence
    \item Financial and legal consequences
    \item Physical harm from long-term use
    \item Carer burnout and stress-related illness in family members
\end{itemize}

\section*{Prognosis and Outcomes}
Recovery is achievable, and family support significantly improves outcomes. Many people reduce or stop use with treatment, though the journey often includes relapses, which are part of recovery rather than failure. Engagement with treatment, family involvement, and management of coexisting mental health conditions are the strongest predictors of success. Families who seek their own support cope better and are better placed to help. With patience and persistence, many people recover fully and rebuild their lives.

\section*{Prevention and Self-Care}
\begin{itemize}
    \item Learn about the substance and its effects
    \item Call a helpline such as the Alcohol and Drug Information Service for advice
    \item Seek family counselling through drug and alcohol services
    \item Set and maintain clear, kind boundaries
    \item Do not provide money or transport that enables use
    \item Encourage healthy routines and interests
    \item Look after your own health, sleep, and support networks
    \item Be patient: change takes time and relapse is common
\end{itemize}

\section*{Sources and Further Reading}
The following reputable sources provide further information for healthcare students and patients seeking reliable, current advice about helping someone with substance misuse.

\begin{itemize}
    \item Alcohol and Drug Foundation. \textit{Supporting someone with a drug or alcohol problem.} https://adf.org.au/
    \item Healthdirect. \textit{Drug and alcohol support services.} https://www.healthdirect.gov.au/drugs
    \item Family Drug Support Australia. \textit{Support for families affected by substance use.} https://www.fds.org.au/
    \item National Alcohol and Other Drug Hotline. \textit{24/7 phone and online support.} https://adf.org.au/
    \item ReachOut and other youth services. \textit{Support for young people and families.} https://au.reachout.com/
    \item SAMHSA and international resources. \textit{Family support for addiction.} https://www.samhsa.gov/
\end{itemize}
